\subsection{Die Minimalität des Anfügungsmodells}

Die Wortmenge wurde als kleinste unter der Erzeugungsregel abgeschlossene
Familie konstruiert. Nach
\FormulaRefAuto{EarlyWordSuccessorGenerationEquation}[thm] führt
\(\sigma_A\) genau diese Regel aus. Daher gilt die bereits bewiesene
Minimalität auch in der Sprache der Anfügungsfunktion.

\noindent\begin{minipage}{\linewidth}
\FormulaThmDeltaK[W3: Minimalität des konkreten Wortmodells]{%
  \begin{aligned}[t]
    &U\subseteq\FiniteWords A\dsep\EmptyWord A\in U\dsep{}\\[-2pt]
    &\forall u\in U\,\forall a\in A\,(\sigma_A(u,a)\in U)
      \vdash U=\FiniteWords A
  \end{aligned}
}{EarlyWordModelMinimality}{%
  \DeltaRow{Mengen}{A\dsep U\dsep u\dsep a}
}
\end{minipage}\par
\begin{tabproofwide}
  \proofstepwidestar[1]{U\subseteq\FiniteWords A}{\rA}
  \proofstepwidestar[2]{\EmptyWord A\in U}{\rA}
  \proofstepwidestar[3]{\forall u\in U\,\forall a\in A\,(\sigma_A(u,a)\in U)}{\rA}
  \proofstepwidestar[]{\FiniteWords A\subseteq\FiniteOccurrenceSets A}{\FormulaRefAuto{FiniteWordSetSubsetUniverse}[thm]}
  \proofstepwidestar[1]{U\subseteq\FiniteOccurrenceSets A}{\FormulaRefAuto{A\subseteq B\dsep B\subseteq C\vdash A\subseteq C}[thm]{1,4}}
  \proofstepwidestar[]{\EmptyWord A=\varnothing}{\FormulaRefAuto{EmptyWordDef}[def]}
  \proofstepwidestar[2]{\varnothing\in U}{\rIE{6,2}}
  \proofstepwidestar[8]{u\in U}{\rA}
  \proofstepwidestar[9]{a\in A}{\rA}
  \proofstepwidestar[1,8]{u\in\FiniteWords A}{\FormulaRefAuto{A\subseteq B\dsep x\in A\vdash x\in B}[thm]{1,8}}
  \proofstepwidestar[3,8]{\forall a\in A\,(\sigma_A(u,a)\in U)}{\rRE{8,\rUE{3}}}
  \proofstepwidestar[3,8,9]{\sigma_A(u,a)\in U}{\rRE{9,\rUE{11}}}
  \proofstepwidestar[1,8,9]{\sigma_A(u,a)=\OccurrenceAdjunction u a}{\FormulaRefAuto{EarlyWordSuccessorGenerationEquation}[thm]{10,9}}
  \proofstepwidestar[1,3,8,9]{\OccurrenceAdjunction u a\in U}{\rIE{13,12}}
  \proofstepwidestar[1,3]{\forall u\in U\,\forall a\in A\,(\OccurrenceAdjunction u a\in U)}{\rUI{\rRI{8,\rUI{\rRI{9,14}}}}}
  \proofstepwidestar[1,2,3]{\FiniteWords A\subseteq U}{\FormulaRefAuto{GeneratedWordMinimality}[thm]{5,7,15}}
  \proofstepwidestar[1,2,3]{U=\FiniteWords A}{\FormulaRefAuto{A\subseteq B\dsep B\subseteq A\vdash A=B}[thm]{1,16}}
\end{tabproofwide}

Damit sind die vier konkreten Grundgesetze bewiesen. Wir halten nun
genau diese Gesetze als Axiome einer abstrakten Wortstruktur fest
und leiten daraus die allgemeine Induktion und Rekursion her.
